\chapter*{\lblIntroduction}
\phantomsection
\addcontentsline{toc}{chapter}{\lblIntroduction}
\markboth{\lblIntroduction}{\lblIntroduction}

\section*{Who this book is for}

You have shipped services in C\#. You know what a connection pool is for, why
idempotency matters on a retried request, what a cancellation token does and
roughly what you would look at first when a service degrades under load. You
have been asked \dash{} by a job advertisement, by a team you joined, by the
fact that the interesting tooling in your field is now written in it \dash{}
to do the same in Python. And you have discovered that a fortnight is enough to
be productive and not nearly enough to stop writing Python that a Python
reviewer quietly rewrites.

This book is for that fortnight and the year after it. It assumes you are a
senior engineer and it does not assume you know Python. It is a translation
dictionary from C\# to Python, and a field manual for putting Python into
production written by somebody who does that for a living.

\section*{What this book is not}

\begin{itemize}[leftmargin=1.4em]
  \item It is not a Python course from zero. It never explains what a loop is,
        and it spends its pages on the places where a decade of C\# produces
        Python that works and reads wrong.
  \item It is not a catalogue of libraries. It picks one tool for each job
        \dash{} \pkg{uv}, \pkg{ruff}, \pkg{pydantic}, \pkg{FastAPI},
        \pkg{SQLAlchemy}, \pkg{pytest} \dash{} pins the version, and says why.
        Alternatives get one sentence.
  \item It is not another book about AI. Part~V is short, and it stops at one
        model call and one structured output.
  \item It is not the LangChain book. That book, \emph{LangChain, LangGraph and
        Async Python}, is already written for an engineer arriving from .NET,
        and this one is its volume zero: the language, the toolchain, the
        idioms and the production habits that book assumes. Agent frameworks
        stay there. Microsoft's Agent Framework stays in the third volume.
        Together the three are meant to be read as one course, and none of
        them repeats a section of another.
\end{itemize}

\mermaidfig{book-map}{Where this book sits: volume zero of a trilogy, with the same reader on both stacks}{the trilogy}

\section*{Why now}

Three things changed in the two years before this book was written, and each
removed an excuse.

Python~\pyver{} made the free-threaded build an officially supported way to
run the interpreter rather than an experiment, which means the one sentence
every .NET engineer knows about Python \dash{} \enquote{it has a global
interpreter lock} \dash{} now needs a second sentence, and
Chapter~\ref{ch:runtime} supplies it and measures it.

The toolchain converged. For most of Python's history the honest answer to
\enquote{what is the equivalent of \code{dotnet restore}} was a shrug and a
list; the answer is now \pkg{uv}, and Chapter~\ref{ch:packaging} is written as
if that were settled, because it is.

And the typing story matured to the point where a C\# engineer's instincts are
an asset rather than a liability, provided they learn exactly where the checker
stops. That is Chapter~\ref{ch:typing}, and it is the chapter this book was
started for.

\section*{How the book is organised}

Five parts, fourteen chapters, each under three thousand words.

\textbf{Part~I \dash{} \lblPartI} (Chapters~\ref{ch:runtime}
to~\ref{ch:typing}) is where .NET engineers fall over first: the runtime, the
environment, the packaging, the typing. Nothing in it is hard and all of it is
different, and getting it wrong is what produces the six-month-old project
that nobody can install.

\textbf{Part~II \dash{} \lblPartII} (Chapters~\ref{ch:objects}
to~\ref{ch:imports}) is the dictionary proper: objects, functions, errors,
imports. Each chapter is a mapping table and the reasons the mapping is not
one to one.

\textbf{Part~III \dash{} \lblPartIII} (Chapter~\ref{ch:asyncio}) is one
chapter, deliberately, because the LangChain book already spends two on the
event loop for an agent workload. This one is the translation from the Task
Parallel Library, and it is the chapter .NET engineers think they do not need.

\textbf{Part~IV \dash{} \lblPartIV} (Chapters~\ref{ch:services}
to~\ref{ch:observability}) is production: a service, a database, a test suite
and the operational surface \dash{} logs, traces, containers, audits \dash{}
each mapped from the ASP.NET Core and EF Core equivalents.

\textbf{Part~V \dash{} \lblPartV} (Chapters~\ref{ch:aitoolkit}
and~\ref{ch:traceassert}) is the AI engineer's kit \dash{} the three things
every SDK in this field is made of \dash{} and the completion of the guiding
project.

Chapters are independent enough to read out of order after Part~I, with two
exceptions: Chapter~\ref{ch:testing} assumes Chapter~\ref{ch:functions}, and
Chapter~\ref{ch:traceassert} assumes everything.

\section*{The guiding project}

The book builds one thing across its length: \textbf{\code{trace-assert}}, a
port to Python of the first layer of \code{agent-eval-bench} \dash{}
deterministic assertions over the execution trace of an agent, written as
\pkg{pytest} tests. \emph{Did the agent call the tool it was supposed to? In
the order it was supposed to? With arguments that validate? Without calling
anything it was forbidden?} Those are questions a test can answer with no
model, no budget and no judgement, and they are the cheapest evaluation an
agent can have.

It arrives in three stages. Chapter~\ref{ch:testing} lays the skeleton \dash{}
the trace model, a fixture and the first two assertions. Chapter~\ref{ch:aitoolkit}
records a real model call as trace events. Chapter~\ref{ch:traceassert}
completes the assertion set, packages it, and closes on a table comparing what
the same instrument cost to build in .NET, in TypeScript and in Python.

The project is not a toy chosen for teaching. It is a thing the author needed,
and the reader finishes the book with a package, a post and an interview
answer, which is the honest reason to choose a project at all.

\section*{The rules this book holds itself to}

\begin{note}
\textbf{Every listing is a file, and every file is run.} The code lives under
\code{code/} as a project of its own, locked to the versions on the title
page, and the repository's build runs every listing and every exercise test on
every change. A listing that has not been through that says so in a box.
\end{note}

\begin{note}
\textbf{Measurements over assertions.} A claim about what is faster, cheaper
or safer comes with a script and a number, or it is labelled as judgement.
Eight measurements are specified across the chapters, all of them free and
finishing on a laptop; until the one behind a claim has run, the claim says
so.
\end{note}

\begin{note}
\textbf{Hard frames.} Fourteen chapters, none over three thousand words, each
publishable on its own, in three releases. The frames exist because the
author's other projects have taught him what an unbounded book becomes, and
the build enforces them rather than trusting anybody to remember.
\end{note}
